% NS-022 prose preparation. No final outcome is asserted in this fragment.
\section{Introduction}

A coding supervisor must decide what a proposed edit changes, what evidence
applies to the resulting program, and who may accept the transition. These
decisions become harder when work spans multiple attempts and interruptions.
A stored passing result can concern an earlier source revision; a completed
task can record an operational decision without establishing a behavioral
property. Our starting point is to make the subject, scope, and authority of
each observation explicit.

We describe a supervision architecture with three separate responsibilities:
semantic state records source-linked facts and unresolved obligations;
operational state records work and decisions; and artifact storage retains
candidate bytes and evidence. A proposed transition must identify the evidence
that supports its declared acceptance conditions. Unknown dependencies,
unsupported translations, and unavailable checks remain visible obligations.
Neither persistence nor a cryptographic commitment upgrades an observation
into a proof of an unstated property.

The paper contributes a design for these evidence boundaries, conditional
arguments for acceptance and reuse, and a frozen protocol for a bounded
repository-repair comparison. The empirical question is whether adding native
source-linked semantic context to matched public raw context changes useful
cold repair under the specified resource policy. Its scope is eight recruited
repository families, two arms, and two nested repetitions per arm. The retained
comparison does not isolate every component of the supervision architecture.
Symbolic routing, warm-cache reuse, publication efficacy, and native proving
performance require separate evidence.

Useful completion requires a candidate that passes the actual nonempty public
and hidden checks within the measured resource rules. Provider failures,
timeouts, and incomplete validation remain in the fixed outcome denominator.
The resulting account therefore separates a successful repair from a returned
proposal, and a measured cost from a configured limit. This distinction also
governs how earlier pilot results, setup work, and an interrupted runtime are
reported alongside the final comparison.
